\documentclass[10pt,a4paper]{amsart}
\usepackage[margin=19mm]{geometry}
\usepackage{amsmath,amssymb,mathtools}
\usepackage[scr=rsfs,cal=euler]{mathalfa}
\usepackage{xfrac}
\usepackage[unicode,hidelinks,bookmarks=false]{hyperref}
\newcommand{\ie}{i.e.\ }
\newcommand{\cf}{cf.\ }
\newcommand{\resp}{respectively\ }
\newcommand{\Subsection}{\subsection}
\providecommand{\nobreakdash}{\nobreak\hbox{-}}
\setlength{\emergencystretch}{2em}
\multlinegap0pt
\usepackage{bm}
\usepackage{amssymb}
\newtheorem{proposition}{Proposition}
\newtheorem{corollary}{Corollary}
\newcommand{\T}{\mathbb{T}}
\newcommand{\X}{\bm{X}}
\newcommand{\abs}[1]{\left\lvert #1 \right\rvert}
\newcommand{\be}{\bm{e}}
\newcommand{\bx}{\bm{x}}
\newcommand{\mc}[1]{\mathcal{#1}}
\newcommand{\norm}[1]{\left\lVert #1 \right\rVert}
\newcommand{\p}{\partial}
\title{The slender body free boundary problem}
\author{Laurel Ohm}
\date{Source: arXiv:2509.16800v1 (CC BY 4.0)}
\begin{document}
\maketitle

\noindent\emph{Source and licence.} The slender body free boundary problem. Version arXiv:2509.16800v1 (CC BY 4.0), distributed by arXiv under the Creative Commons Attribution 4.0 International licence, http://creativecommons.org/licenses/by/4.0/, as recorded in the arXiv licence metadata for this exact version and shown on the abstract page https://arxiv.org/abs/2509.16800v1. Full licence notice: \texttt{licenses/arxiv-2509.16800-CC-BY-4.0.txt}.\par\smallskip

\noindent\textit{Editorial scope.} Counted unit: the article's corollary cor:TDP\_solve (source0.tex lines 1226--1236). The article's complete original proof of it is delivered with it (source0.tex lines 1239--1275, one \texttt{proof} environment of 37 lines). No mathematics is written, repaired or re-derived: the statement and the proof are the article's own lines, in the article's own notation, and no retained line is substituted or re-flowed.  Retained uncounted as the source-local context the proof needs: lines 142--144; lines 151--153; lines 349--351; lines 809--845; lines 1084--1086; lines 1096--1107.  Nothing was omitted from any retained span, so the package carries no omission record.  No retained line contains a citation, so the wrapper carries no bibliography and no bibliography entry.  No retained line contains a figure environment, an image-inclusion command or drawing code, so no image is delivered and the package declares no asset dependency.  Editorial formatting only, in addition: an AMS \texttt{amsart} preamble that restates the article's own declarations and macros used by the retained lines, namely the environments \texttt{proposition}, \texttt{corollary} on separate counters, together with the standard packages those lines need.  The retained environments print in this document as Proposition 1, Corollary 1 in order of appearance, whereas the article prints the same items with the numbering of its own full document.  The complete native source of the article as distributed in the arXiv e-print for this exact version is bundled unchanged as \texttt{sources/185412/source0.tex}.
\begin{equation}\label{eq:SigmaEps}
  \Sigma_\epsilon(t) = \big\{ \bx\in \X_s(s,t)^\perp\,:\, {\rm dist}(\bx,\X(s,t))<\epsilon\,, \, s\in \T \big\}\,.
\end{equation}

\begin{equation}\label{eq:rstar}
  0<r_\star = r_\star(c_\Gamma,\norm{\X_{ss}}_{L^\infty(\T)}) \le \min\bigg\{\frac{c_\Gamma}{2}, \frac{1}{2\norm{\X_{ss}}_{L^\infty(\T)}}\bigg\}
\end{equation}

\begin{equation}\label{eq:TDP}
\mc{Q}_\epsilon[\tau] = -\big(\mc{L}_\epsilon[\bm{g}]\big)_s\cdot\X_s\,.
\end{equation}

\begin{proposition}\label{prop:TDP_reform}
The tension determination problem \eqref{eq:TDP} may be reformulated as 
  \begin{equation}\label{eq:TDP_reform}
    (I +\mc{K})[\tau] = (\mc{G}_0+\mc{G}_\epsilon+\mc{G}_+')[\bm{g}]
  \end{equation}
where, for any $0<\alpha<\gamma<1$, the operator $\mc{K}$ satisfies
\begin{equation}
  \norm{\mc{K}[\tau]}_{C^{1,\gamma}(\T)} \le c(\epsilon,\norm{\X}_{C^{3,\alpha}},c_\Gamma)\norm{\tau}_{C^{1,\alpha}(\T)}\,.
\end{equation}
Given two nearby filaments with centerlines $\X^{(a)}$, $\X^{(b)}$ in $C^{3,\alpha}(\T)$, the difference between the corresponding operators $\mc{K}^{(a)}$ and $\mc{K}^{(b)}$ satisfies
\begin{equation}
  \norm{(\mc{K}^{(a)}-\mc{K}^{(b)})[\tau]}_{C^{1,\gamma}(\T)} \le c(\epsilon,\|\X^{(a)}\|_{C^{3,\alpha}},\|\X^{(b)}\|_{C^{3,\alpha}},c_\Gamma)\norm{\X^{(a)}-\X^{(b)}}_{C^{3,\alpha}}\norm{\tau}_{C^{1,\alpha}(\T)}\,.
\end{equation}

Meanwhile, the main term on the right hand side is given explicitly by
\begin{equation}
  \mc{G}_0[\bm{g}] = (I-\overline{\mc{L}}_\epsilon^{\rm tang}\p_{ss})^{-1}\p_s\overline{\mc{L}}_\epsilon^{\rm tang}\bigg[\be_{\rm t}\cdot\bm{g}-\int_\T\be_{\rm t}\cdot\bm{g}\,ds \bigg]\,,
\end{equation}
while the remainder terms satisfy 
\begin{equation}
\begin{aligned}
  \norm{\mc{G}_\epsilon[\bm{g}]}_{C^{1,\alpha}}&\le c(\|\X\|_{C^{2,\alpha^+}},c_\Gamma)\,\epsilon^{1-\alpha^+}\abs{\log\epsilon}^3\norm{\bm{g}}_{C^{0,\alpha}}\\
  %
  \norm{\mc{G}_+'[\bm{g}]}_{C^{1,\gamma}}&\le c(\epsilon,\norm{\X}_{C^{3,\alpha}},c_\Gamma)\norm{\bm{g}}_{C^{0,\alpha}}
\end{aligned}
\end{equation}
along with the Lipschitz bounds
\begin{equation}
\begin{aligned}
  &\norm{(\mc{G}_\epsilon^{(a)}-\mc{G}_\epsilon^{(b)})[\bm{g}]}_{C^{1,\alpha}}\\
  &\quad\qquad\le c(\|\X^{(a)}\|_{C^{2,\alpha^+}},\|\X^{(b)}\|_{C^{2,\alpha^+}},c_\Gamma)\,\epsilon^{1-\alpha^+}\abs{\log\epsilon}^3\norm{\X^{(a)}-\X^{(b)}}_{C^{2,\alpha^+}}\norm{\bm{g}}_{C^{0,\alpha}}\\
  %
  &\norm{(\mc{G}_+^{'(a)}-\mc{G}_+^{'(b)})[\bm{g}]}_{C^{1,\gamma}}\le c(\epsilon,\|\X^{(a)}\|_{C^{3,\alpha}},\|\X^{(b)}\|_{C^{3,\alpha}},c_\Gamma)\norm{\X^{(a)}-\X^{(b)}}_{C^{2,\gamma^+}}\norm{\bm{g}}_{C^{0,\alpha}}
\end{aligned}
\end{equation}
for any $\alpha^+>\alpha$ and $\gamma^+>\gamma$.
\end{proposition}

\begin{equation}\label{eq:TDP_h}
(I+\mc{K})[\psi] = h(s)
\end{equation}

\begin{proposition}\label{prop:tdp_qual}
  Given $\Sigma_\epsilon$ as in \eqref{eq:SigmaEps}-\eqref{eq:rstar} with centerline $\X\in C^{3,\alpha}(\T)$, $0<\alpha<1$, and given $h\in C^{1,\alpha}(\T)$, there exists a unique solution $\psi\in C^{1,\alpha}(\T)$ to the equation \eqref{eq:TDP_h} which satisfies
  \begin{equation}\label{eq:tauh_bd}
  \norm{\psi}_{C^{1,\alpha}(\T)}\le c(\epsilon,\norm{\X}_{C^{3,\alpha}},c_\Gamma)\norm{h}_{C^{1,\alpha}(\T)}\,.
  \end{equation}

Moreover, given nearby filaments with centerlines $\X^{(a)}$, $\X^{(b)}$ in $C^{3,\alpha}(\T)$, the difference $\psi^{(a)}-\psi^{(b)}= \big((I+\mc{K}^{(a)})^{-1}-(I+\mc{K}^{(b)})^{-1}\big)[h]$ satisfies  
  \begin{equation}\label{eq:tauh_lip}
  \norm{\psi^{(a)}-\psi^{(b)}}_{C^{1,\alpha}(\T)}\le c(\epsilon,\|\X^{(a)}\|_{C^{3,\alpha}},\|\X^{(b)}\|_{C^{3,\alpha}},c_\Gamma)\norm{\X^{(a)}-\X^{(b)}}_{C^{3,\alpha}}\norm{h}_{C^{1,\alpha}(\T)}
  \end{equation}
  for $\alpha^+>\alpha$.
\end{proposition}

\begin{corollary}[Solvability of tension determination problem]\label{cor:TDP_solve}
Let $\X\in C^{3,\alpha}(\T)$, $0<\alpha<1$, be the centerline of a filament $\Sigma_\epsilon$ as in \eqref{eq:SigmaEps}-\eqref{eq:rstar}. Given $\bm{g}\in C^{0,\alpha}(\T)$, the reformulated tension determination problem of Proposition \ref{prop:TDP_reform} admits a unique solution $\tau\in C^{1,\alpha}(\T)$ satisfying 
  \begin{equation}
    \norm{\tau}_{C^{1,\alpha}(\T)} \le c(\epsilon,\norm{\X}_{C^{3,\alpha}},c_\Gamma)\norm{\bm{g}}_{C^{0,\alpha}}\,.
  \end{equation}
%
Moreover, given nearby filaments with centerlines $\X^{(a)}$, $\X^{(b)}$ in $C^{3,\alpha}(\T)$, the difference between the solutions to the corresponding tension determination problems satisfies 
  \begin{equation}
    \norm{\tau^{(a)}-\tau^{(b)}}_{C^{1,\alpha}(\T)} \le c(\epsilon,\|\X^{(a)}\|_{C^{3,\alpha}},\|\X^{(b)}\|_{C^{3,\alpha}},c_\Gamma)\norm{\X^{(a)}-\X^{(b)}}_{C^{3,\alpha}(\T)}\norm{\bm{g}}_{C^{0,\alpha}(\T)}\,.
  \end{equation}  
\end{corollary}

\begin{proof}
Given the reformulation of the tension determination problem in Proposition \ref{prop:TDP_reform}, we may write 
\begin{equation}
  \tau = (I+\mc{K})^{-1}\big(\mc{G}_0+ \mc{G}_\epsilon + \mc{G}_+'\big)[\bm{g}] \,.
\end{equation}
%
Using Proposition \ref{prop:tdp_qual}, we may then estimate
\begin{equation}
\begin{aligned}
  \norm{\tau}_{C^{1,\alpha}}&\le c(\epsilon,\norm{\X}_{C^{3,\alpha}},c_\Gamma)\norm{(\mc{G}_0+ \mc{G}_\epsilon + \mc{G}_+')[\bm{g}]}_{C^{1,\alpha}}\\
  %
  &\le c(\epsilon,\norm{\X}_{C^{3,\alpha}},c_\Gamma)\norm{\bm{g}}_{C^{0,\alpha}}\,,
\end{aligned}
\end{equation}
where we have used the coarsest bounds for $\bm{g}$ in Proposition \ref{prop:TDP_reform}. 


Furthermore, given two nearby filaments with centerlines $\X^{(a)}$ and $\X^{(b)}$, we may bound
\begin{equation}
\begin{aligned}
  \norm{\tau^{(a)}-\tau^{(b)}}_{C^{1,\alpha}} &\le 
  \norm{ \big((I+\mc{K}^{(a)})^{-1}-(I+\mc{K}^{(b)})^{-1}\big)\big(\mc{G}_0^{(a)}+ \mc{G}_\epsilon^{(a)} + \mc{G}_+^{'(a)}\big)[\bm{g}] }_{C^{1,\alpha}}\\
  &\quad + \norm{ (I+\mc{K}^{(b)})^{-1}\big(\mc{G}_0^{(a)}-\mc{G}_0^{(b)}+ \mc{G}_\epsilon^{(a)}-\mc{G}_\epsilon^{(b)} + \mc{G}_+^{'(a)}-\mc{G}_+^{'(b)}\big)[\bm{g}] }_{C^{1,\alpha}}\\
  %
  %
  &\hspace{-2cm}\le c(\epsilon,\|\X^{(a)}\|_{C^{3,\alpha}},\|\X^{(b)}\|_{C^{3,\alpha}},c_\Gamma)\norm{\X^{(a)}-\X^{(b)}}_{C^{3,\alpha}}\norm{\big(\mc{G}_0^{(a)}+ \mc{G}_\epsilon^{(a)} + \mc{G}_+^{'(a)}\big)[\bm{g}] }_{C^{1,\alpha}}\\
  &\hspace{-2cm}\quad + c(\epsilon,\norm{\X^{(b)}}_{C^{3,\alpha}},c_\Gamma)\norm{\big(\mc{G}_0^{(a)}-\mc{G}_0^{(b)}+ \mc{G}_\epsilon^{(a)}-\mc{G}_\epsilon^{(b)} + \mc{G}_+^{'(a)}-\mc{G}_+^{'(b)}\big)[\bm{g}] }_{C^{1,\alpha}}\\
  %
  %
  &\hspace{-2cm}\le c(\epsilon,\|\X^{(a)}\|_{C^{3,\alpha}},\|\X^{(b)}\|_{C^{3,\alpha}},c_\Gamma)\norm{\X^{(a)}-\X^{(b)}}_{C^{3,\alpha}}\norm{\bm{g}}_{C^{0,\alpha}}\,,
\end{aligned}
\end{equation}
where we have used Propositions \ref{prop:TDP_reform} and \ref{prop:tdp_qual}.
%
%
%
\end{proof}

\end{document}
